\documentclass[a4paper,14pt]{extarticle}
\usepackage[a4paper,left=25mm,right=15mm,top=20mm,bottom=20mm]{geometry}
\usepackage{fontspec}
\usepackage{polyglossia}
\setdefaultlanguage{russian}
\setotherlanguage{english}
\setmainfont{Times New Roman}
\usepackage{amsmath,amssymb}
\usepackage{booktabs}
\usepackage{array}
\usepackage{caption}
\usepackage{graphicx}
\usepackage{indentfirst}
\usepackage{setspace}
\onehalfspacing
\setlength{\parindent}{1.25cm}

\captionsetup{labelformat=empty,labelsep=none,justification=centering,singlelinecheck=false,font=normalfont}

\begin{document}

\begin{center}
Домашнее задание: отрывки из ВКР
\end{center}

\noindent Тема ВКР: «Гибридная нейросетевая модель регионального прогноза погоды на основе графовых нейронных сетей с усвоением данных и статистическим постпроцессингом».

\noindent Студент: Табаков А.~С. \quad Научный руководитель: Пененко~А.~В.

\bigskip

\section*{1. Отрывок с таблицей}

Для оценки эффективности усвоения данных на пилотной модели НСКО (64\,$\times$\,32, 15~переменных, область 53--57$^{\circ}$~с.\,ш., 79--85$^{\circ}$~в.\,д.) проведена серия экспериментов с двумя алгоритмами~--- Nudging и оптимальной интерполяцией (ОИ)~--- при доле наблюдений 10\,\% и стандартном отклонении ошибки наблюдений $\sigma_o = 0.5$. Сравнение базового прогноза без усвоения, прогноза с Nudging и прогноза с ОИ при оптимальных параметрах приведено в таблице~1.

\begin{table}[h!]
\caption{\hfill Таблица~1}
\vspace{-0.3em}
\begin{center}
Сравнение методов усвоения данных на пилотной модели НСКО (200~сэмплов, +6~ч)
\end{center}
\centering\small
\begin{tabular}{lcccc}
\toprule
Метод & Параметры & Region Skill & RMSE $t_{2m}$, $^{\circ}$C & ACC \\
\midrule
Без DA   & ---                                       & 0.629 & 1.80 & 0.451 \\
Nudging  & $\alpha = 0.9$                             & 0.602 & 1.65 & 0.512 \\
ОИ       & $L = 100$~км, $\sigma_o = 0.5$             & 0.414 & 0.62 & 0.704 \\
\bottomrule
\end{tabular}
\end{table}

Как видно из таблицы~1, оптимальная интерполяция значительно превосходит Nudging: RMSE по приземной температуре снижается с 1.80~$^{\circ}$C до 0.62~$^{\circ}$C (улучшение в 2.9~раза), а аномальная корреляция (ACC) возрастает с 0.451 до 0.704. Это объясняется тем, что ОИ учитывает пространственную корреляционную структуру ошибок прогноза, тогда как Nudging применяет коррекцию точечно. Полученные результаты послужили основанием для перехода к экспериментам на мультирезолюционной модели, описываемым в следующем разделе.

\bigskip

\section*{2. Отрывок с рисунком}

Для качественной оценки эффекта оптимальной интерполяции на пилотной модели НСКО проведено сравнение полей ошибок базового прогноза и прогноза с применением ОИ относительно поля «истины» из реанализа ERA5. На рисунке~1 показаны: эталонное температурное поле, поле ошибки базового прогноза и поле ошибки итоговой модели с ОИ.

\begin{figure}[h!]
\centering
\includegraphics[width=0.95\linewidth]{../slides_visuals/FINAL_RESULT.png}
\caption{Рисунок~1~--- Сравнение полей ошибок базового прогноза и прогноза с оптимальной интерполяцией относительно ERA5}
\end{figure}

Как видно на рисунке~1, RMSE приземной температуры по региону снижается с 2.21~$^{\circ}$C для базового прогноза до 1.37~$^{\circ}$C при использовании оптимальной интерполяции, то есть на 38\,\%. Поле ошибки заметно «осветляется» --- крупные положительные смещения в южной части области и отрицательные в северо-восточной существенно уменьшаются по амплитуде. Полученный результат подтверждает, что ОИ корректирует не отдельные точки, а пространственную структуру ошибки в целом, что согласуется с количественными оценками раздела~7.6 ВКР.

\bigskip

\section*{3. Отрывок с формулой}

Алгоритм оптимальной интерполяции реализует статистически оптимальную взвешенную комбинацию фонового поля прогноза и наблюдений с учётом пространственной корреляционной структуры ошибок. Анализное состояние вычисляется по классической формуле Гандина:

\[
x_a \;=\; x_b \;+\; K\,(y_o - H x_b), \qquad K \;=\; B H^{T}\bigl(H B H^{T} + R\bigr)^{-1},
\]

где $x_b$~--- фоновое поле (прогноз модели), $y_o$~--- вектор наблюдений, $H$~--- оператор наблюдений, $B$ и $R$~--- матрицы ковариаций ошибок прогноза и наблюдений соответственно, $K$~--- матрица усиления Калмана. Матрица $B$ в работе задаётся гауссовой моделью пространственных корреляций

\[
B_{ij} \;=\; \sigma_b^{2}\,\exp\!\left(-\dfrac{d_{ij}^{2}}{L^{2}}\right),
\]

где $L$~--- радиус корреляции, $d_{ij}$~--- расстояние между узлами. Оптимальные значения параметров ($L = 100$~км, $\sigma_o = 0.5$) подобраны в ходе серии из 176~конфигураций, описанной в разделе~7.6 ВКР.

\end{document}
